\section{Related Work}
\label{sec:related}

\fan{for each point/paragraph. We should have one sentence summarizing our difference/novelty over existing related work. or improve the last concluding paragraph to summarize. }
\yifan{Resolved: added distinguishing sentence at end of each paragraph + comparison table.}

\paragraph{Diffusion Language Models.}
Masked diffusion language models corrupt text by replacing tokens with \texttt{[MASK]} and train a model to reverse the process~\citep{austin2021d3pm, sahoo2024mdlm, lou2024sedd}.
LLaDA~\citep{nie2025llada} scaled this paradigm to 8B parameters, LLaDA~2.0~\citep{llada2} to 100B via mixture-of-experts, and LLaDA~2.1~\citep{llada21} introduced token editing with confidence-based decoding.
DREAM~\citep{dream2025} introduced the logit shift technique, aligning the diffusion objective with the AR model's $\text{logits}[i] \to \text{token}[i{+}1]$ mapping.
Block Diffusion~\citep{blockdiffusion2025} generates fixed-size blocks autoregressively while denoising tokens within each block.

Converting pretrained AR models into diffusion models is a more data-efficient alternative.
\citet{shi2024scaling} showed that fine-tuning with a masked diffusion objective reduces training cost significantly.
Self-Distillation Through Time~\citep{hou2024sdtt} and Efficient-DLM~\citep{efficientdlm2025} further streamline the pipeline.
SDAR~\citep{sdar2025} converts AR models via full-model training on ${\sim}$50B tokens for block-parallel generation, and NBDiff~\citep{nbdiff2025} extends this with causal prefix constraints.
TiDAR~\citep{tidar2025} proposes a sequence-level hybrid that drafts via diffusion and verifies autoregressively.
Unlike these approaches, Strided DLLM combines strict causal attention with logit-shifted prediction from MASK positions---by maximally respecting the pretrained AR model's attention and prediction patterns, our conversion is far more data-efficient and closes the quality gap to the base AR model.

% \subsection{Parallel Decoding for LLMs}
% \label{sec:rw_parallel}

\paragraph{Speculative decoding and multi-token prediction for LLMs.}
Speculative decoding~\citep{leviathan2023speculative, chen2023speculative} accelerates AR models by drafting tokens with a fast model and verifying them in parallel, provably preserving the target distribution.
Extensions include Medusa~\citep{medusa2024} (multiple prediction heads), the EAGLE family~\citep{eagle2024, eagle3} (feature-level drafting), and SpecInfer~\citep{specinfer2023} (tree-structured verification).
% \paragraph{ and register-based AR.}
Multi-token prediction (MTP) trains models to predict multiple future tokens simultaneously~\citep{gloeckle2024mtp}.
% ; DeepSeek-V3~\citep{deepseekv3} uses MTP in production.
\citet{applemtp2025} proposed gated sparse expert LoRA for MTP---the closest prior work to our conditional LoRA, though designed for multi-head prediction rather than diffusion verification.
Consistency LLMs~\citep{cllms2024} and Lookahead Decoding~\citep{lookahead2024} achieve acceleration via Jacobi iteration.
ISD differs from these by using a single model with training-time stride capability to both achieve quality guarantee and deliver high latency and throughput gains under high concurrency.

\paragraph{DLLM-specific decoding.}
FastDLLM~\citep{fastdllm2025} enables confidence-aware parallel decoding with KV cache reuse.
Jacobi Forcing~\citep{jacobiforcing2025} distills diffusion models for fewer-step convergence but degrades at larger batch sizes.
SpecDiff~\citep{specdiff2024} uses a diffusion model as a drafter for AR verification, while Spiffy~\citep{spiffy2025} and Free Draft-and-Verification~\citep{freedraftverify2025} explore self-speculative approaches.
WeDLM~\citep{wedlm2025} reconciles diffusion with causal attention for KV cache reuse.
These methods rely on confidence-based acceptance or iterative denoising, which either lacks formal quality guarantees or incurs high compute overhead; our ISD provides provable AR-quality output in a single stride--introspection cycle.

\paragraph{Comparison of parallel decoding approaches.}
Table~\ref{tab:comparison} summarizes the key design dimensions across parallel decoding paradigms.

\begin{table*}[t]
\centering
\caption{\textbf{Comparison of parallel decoding approaches.} \ding{51} = yes, \ding{55} = no. OH = overhead. Qual.\ Guar.\ = quality guarantee. Low-C / High-C = latency and throughput under low / high concurrency. \fan{font size too small}}
\label{tab:comparison}
\setlength{\tabcolsep}{2pt}
\resizebox{\textwidth}{!}{
\begin{tabular}{lcccccccccccccc}
\toprule
\textbf{Method} & \textbf{Examples} & \textbf{Arch.\ Change} & \textbf{Logit Shift} & \textbf{Attn Mask} & \textbf{Input} & \textbf{Compute OH} & \textbf{Memory OH} & \textbf{AR Infra} & \textbf{Training} & \textbf{Quality} & \textbf{Qual.\ Guar.} & \textbf{Low-C} & \textbf{High-C} \\
\midrule
MTP & DeepSeek-V3 & Extra heads & \ding{51} & Causal & Clean & Extra heads & Extra heads & \ding{51} & From pretrain & Comparable & --- & Moderate & Good \\
Spec.\ Decoding & EAGLE-3, Medusa & Draft model & \ding{51} & Causal & Clean & Draft overhead & Extra draft head & \ding{51} & Draft model & Lossless & $p/q$ & Good & Degrades \\
Bidir.\ DLLM & LLaDA & None & \ding{55} & Bidir. & Clean+masked & $T{\times}L$ denoise & No KV cache & \ding{55} & From pretrain & Lower & Confidence & Good & Poor \\
Block-causal DLLM & SDAR, NBDiff & None & \ding{55} & Block-causal & Clean+masked & $T{\times}N$ denoise & Minimal & \ding{55} & Conversion & Lower & Confidence & Moderate & Poor \\
\midrule
\textbf{Strided DLLM (Ours)} & --- & \textbf{None} & \ding{51} & \textbf{Causal} & \textbf{Clean+masked} & \textbf{Verify tokens} & \textbf{N/A} & \ding{51} & \textbf{Eff.\ conv.} & \textbf{Comp./Lossless} & $\mathbf{p/q}$ & \textbf{Best} & \textbf{Best} \\
\bottomrule
\end{tabular}
}
\end{table*}

Despite these advances, fundamental gaps remain between DLLMs and AR models in training scalability, inference alignment, compute efficiency, and infrastructure compatibility. We analyze these gaps quantitatively in the next section.
